%! TeX program = pdflatex
%! BIB program = biber
% =========================================================================
% MASTER FILE: BUKU LAPORAN TUGAS AKHIR (SEMHAS & SIDANG AKHIR)
% Departemen Teknologi Kedokteran - Fakultas Kedokteran dan Kesehatan ITS
% Mengacu pada SK Rektor ITS No. 280/IT2/T/HK.00.01/2022 & Standar B201 ITS
% =========================================================================

\newcommand{\doctypename}{TUGAS AKHIR}
\newcommand{\engdoctypename}{FINAL PROJECT}

% Judul dokumen
\title{Buku Tugas Akhir ITS - Teknologi Kedokteran}
\author{\name}

% Margin 4433 tetap pada setiap halaman: atas 4 cm, bawah 3 cm,
% kiri (sisi jilid) 4 cm, kanan 3 cm.
\documentclass[12pt,oneside]{report}
\usepackage[a4paper,top=40mm,bottom=30mm,left=40mm,right=30mm]{geometry}

% Pengaturan ukuran spasi (1 / singlespacing)
\usepackage[singlespacing]{setspace}

% Pengaturan detail pada file PDF
\usepackage[pdfauthor={\@author},bookmarksnumbered,pdfborder={0 0 0}]{hyperref}

% Pengaturan jenis karakter
\usepackage[utf8]{inputenc}
\usepackage[T1]{fontenc}

% Pengaturan pewarnaan
\usepackage[table,xcdraw]{xcolor}

% Pengaturan kutipan artikel
\usepackage[style=apa, backend=biber]{biblatex}

% Package lainnya
\usepackage{changepage}
\usepackage{enumitem}
\usepackage{eso-pic}
\usepackage{txfonts} % Font times
\usepackage{etoolbox}
\usepackage{graphicx}
\usepackage{lipsum}
\usepackage{longtable}
\usepackage{tabularx}
\usepackage{wrapfig}
\usepackage{float}

% Halaman laporan dicetak satu sisi; tidak perlu halaman kosong penyela.

% Nomor halaman di tengah bawah.
\usepackage{fancyhdr}
\fancyhf{}
\renewcommand{\headrulewidth}{0pt}
\pagestyle{fancy}
\fancyfoot[C]{\thepage}
\patchcmd{\chapter}{plain}{fancy}{}{}
\patchcmd{\chapter}{empty}{plain}{}{}

% Command untuk bulan
\newcommand{\MONTH}{%
  \ifcase\the\month
  \or Januari% 1
  \or Februari% 2
  \or Maret% 3
  \or April% 4
  \or Mei% 5
  \or Juni% 6
  \or Juli% 7
  \or Agustus% 8
  \or September% 9
  \or Oktober% 10
  \or November% 11
  \or Desember% 12
  \fi}
\newcommand{\ENGMONTH}{%
  \ifcase\the\month
  \or January% 1
  \or February% 2
  \or March% 3
  \or April% 4
  \or May% 5
  \or June% 6
  \or July% 7
  \or August% 8
  \or September% 9
  \or October% 10
  \or November% 11
  \or December% 12
  \fi}

% Gaya Heading 1--4, Gambar, dan Tabel
% Heading 1--4: bab, subbab, anak subbab, dan subanak subbab.
\usepackage{titlesec}
\titleformat{\chapter}[display]{\bfseries\fontsize{14}{17}\selectfont\centering}{BAB \Roman{chapter}}{0pt}{\MakeUppercase}
\titleformat{\section}{\bfseries\fontsize{13}{16}\selectfont}{\MakeUppercase{\thesection}}{1ex}{}
\titleformat{\subsection}{\bfseries\fontsize{12}{15}\selectfont}{\MakeUppercase{\thesubsection}}{1ex}{}
\titleformat{\subsubsection}{\bfseries\fontsize{12}{15}\selectfont}{\MakeUppercase{\thesubsubsection}}{1ex}{}
\titlespacing{\chapter}{0ex}{0ex}{0ex}
\titlespacing{\section}{0ex}{1ex}{0ex}
\titlespacing{\subsection}{0ex}{0.5ex}{0ex}
\titlespacing{\subsubsection}{0ex}{0.5ex}{0ex}

% Paragraf isi: 12 pt normal, inden baris pertama, rata kiri-kanan, jarak akhir 1ex.
\newcommand{\paragraphtextstyle}{%
  \normalfont\rmfamily\fontsize{12}{15}\selectfont
  \setlength{\parindent}{2em}%
  \setlength{\parskip}{1ex}%
}

% List: rapat antarbutir, jarak 1ex dari paragraf sebelumnya, jarak setelah daftar tetap ada.
\newcommand{\listtextstyle}{%
  \setlist{itemsep=0pt,parsep=0pt,topsep=0pt,partopsep=0pt,%
    after=\vspace{\parskip}}%
}
\listtextstyle

% Isi daftar awal: 11 pt Times-compatible (txfonts), tanpa huruf tebal atau titik penghubung.
\usepackage[titles]{tocloft}
\renewcommand{\cftchapfont}{\normalfont}
\renewcommand{\cftchappagefont}{\normalfont}
\renewcommand{\cftsecfont}{\normalfont}
\renewcommand{\cftsecpagefont}{\normalfont}
\renewcommand{\cftsubsecfont}{\normalfont}
\renewcommand{\cftsubsecpagefont}{\normalfont}
\renewcommand{\cftsubsubsecfont}{\normalfont}
\renewcommand{\cftsubsubsecpagefont}{\normalfont}
\renewcommand{\cftfigfont}{\normalfont}
\renewcommand{\cftfigpagefont}{\normalfont}
\renewcommand{\cfttabfont}{\normalfont}
\renewcommand{\cfttabpagefont}{\normalfont}
\renewcommand{\cftchappresnum}{BAB~}
\setlength{\cftchapnumwidth}{4em}
\setlength{\cftbeforechapskip}{1ex}
\setlength{\cftbeforesecskip}{1ex}
\setlength{\cftbeforesubsecskip}{1ex}
\setlength{\cftbeforesubsubsecskip}{1ex}
\setlength{\cftbeforefigskip}{1ex}
\setlength{\cftbeforetabskip}{1ex}
\renewcommand{\cftchapleader}{\hfill}
\renewcommand{\cftsecleader}{\hfill}
\renewcommand{\cftsubsecleader}{\hfill}
\renewcommand{\cftsubsubsecleader}{\hfill}
\renewcommand{\cftfigleader}{\hfill}
\renewcommand{\cfttableader}{\hfill}

% Gambar dan Tabel: 10 pt reguler, rata tengah, dengan jarak setelah caption.
\renewcommand{\figurename}{Gambar}
\renewcommand{\tablename}{Tabel}
\makeatletter
\setlength{\abovecaptionskip}{3pt}
\setlength{\belowcaptionskip}{1ex}
\long\def\@makecaption#1#2{%
  \vskip\abovecaptionskip
  \begingroup
  \fontsize{10}{12}\selectfont
  \sbox\@tempboxa{#1. #2}%
  \ifdim \wd\@tempboxa >\hsize
    \centering #1. #2\par
  \else
    \global \@minipagefalse
    \hb@xt@\hsize{\hfil\box\@tempboxa\hfil}%
  \fi
  \endgroup
  \vskip\belowcaptionskip}
\makeatother


% =========================================================================
% Variabel proposal dan laporan magang Teknologi Kedokteran ITS
% Fakultas Kedokteran dan Kesehatan
% =========================================================================
% Petunjuk: Sesuaikan variabel di bawah ini dengan data Anda masing-masing.

% -------------------------------------------------------------------------
% 1. IDENTITAS MAHASISWA
% -------------------------------------------------------------------------
\newcommand{\name}{Nama Lengkap Mahasiswa}
\newcommand{\authorname}{Nama Belakang, Nama Depan}
\newcommand{\nickname}{Nama Panggilan}
\newcommand{\nrp}{5049231XXX}

% -------------------------------------------------------------------------
% 2. PEMBIMBING MAGANG
% -------------------------------------------------------------------------
\newcommand{\advisor}{Nama Pembimbing Utama, Gelar}
\newcommand{\advisornip}{xxxxxxxxxxxxxx}

\newcommand{\coadvisorrole}{Dokter Pembimbing} % atau 'Ko-pembimbing' / 'Dosen Pembimbing II'
\newcommand{\engcoadvisorrole}{Co-Advisor}     % atau 'Doctor Advisor' / 'Co-Advisor'
\newcommand{\coadvisor}{Nama Pembimbing Pendamping, Gelar}
\newcommand{\coadvisornip}{NIP/NPP Dosen Pembimbing Pendamping}


% -------------------------------------------------------------------------
% 3. JUDUL MAGANG
% -------------------------------------------------------------------------
\newcommand{\internshiptitle}{[Judul proposal/laporan magang]}

% -------------------------------------------------------------------------
% 4. TEMPAT DAN TAHUN
% -------------------------------------------------------------------------
\newcommand{\place}{Surabaya}
% Data proposal magang; lengkapi sesuai informasi instansi.
\newcommand{\internshipperiod}{Periode pelaksanaan magang}
\newcommand{\internshipcompany}{Nama instansi tempat magang}
\newcommand{\internshipaddress}{Alamat instansi tempat magang}
\newcommand{\internshipphone}{Nomor telepon instansi}
\newcommand{\internshipemail}{Email instansi}
\newcommand{\internshipcoordinator}{Nama Koordinator Magang}
\newcommand{\internshipcoordinatornip}{xxxxxxxxxxxxxx}
\newcommand{\internshipsupervisor}{Nama Pembimbing Lapangan}

% -------------------------------------------------------------------------
% 5. PROGRAM STUDI & DEPARTEMEN (Khusus Teknologi Kedokteran ITS)
% -------------------------------------------------------------------------
\newcommand{\studyprogram}{Teknologi Kedokteran}
\newcommand{\engstudyprogram}{Medical Technology}

\newcommand{\department}{Teknologi Kedokteran}
\newcommand{\engdepartment}{Medical Technology}

% -------------------------------------------------------------------------
% 6. FAKULTAS (Fakultas Kedokteran dan Kesehatan ITS)
% -------------------------------------------------------------------------
\newcommand{\faculty}{Kedokteran dan Kesehatan}
\newcommand{\engfaculty}{Medicine and Health}

\newcommand{\facultyshort}{FKK}
\newcommand{\engfacultyshort}{FMH}

% -------------------------------------------------------------------------

\renewcommand{\coursecode}{\tacoursecode}
\renewcommand{\coverdoctypename}{LAPORAN MAGANG}
\newcommand{\approvaltitlelabel}{Judul Laporan}

% Tambahkan format tanda hubung yang benar di sini
\hyphenation{
  ro-ket
  me-ngem-bang-kan
  per-hi-tu-ngan
  tek-no-lo-gi
  me-la-ku-kan
  ber-so-si-al-i-sa-si
}

% Menambahkan resource daftar pustaka
\addbibresource{pustaka/pustaka.bib}

% Pengaturan format potongan kode
\usepackage{listings}
\definecolor{comment}{RGB}{0,128,0}
\definecolor{string}{RGB}{255,0,0}
\definecolor{keyword}{RGB}{0,0,255}
\lstdefinestyle{codestyle}{
  commentstyle=\color{comment},
  stringstyle=\color{string},
  keywordstyle=\color{keyword},
  basicstyle=\footnotesize\ttfamily,
  numbers=left,
  numberstyle=\tiny,
  numbersep=5pt,
  frame=lines,
  breaklines=true,
  prebreak=\raisebox{0ex}[0ex][0ex]{\ensuremath{\hookleftarrow}},
  showstringspaces=false,
  upquote=true,
  tabsize=2,
}
\lstset{style=codestyle}

\begin{document}

% =========================================================================
% HALAMAN SAMPUL LUAR (Hardcover Biru ITS - Lampiran 6)
% =========================================================================
\newcommand\covercontents{sampul/konten-id.tex}
\AddToShipoutPictureBG*{%
  \AtPageLowerLeft{%
    \color[RGB]{5,91,160}\rule{\paperwidth}{198mm}%
    \hspace{-\paperwidth}%
    \raisebox{206mm}{\color[RGB]{18,173,163}\rule{\paperwidth}{11mm}}%
    \hspace{-\paperwidth}%
    \hspace{46mm}\raisebox{228mm}{\includegraphics[width=40mm]{sampul/gambar/logo-its.png}}%
  }%
}

\thispagestyle{empty}
\newgeometry{top=105mm,left=41mm,right=20mm,bottom=20mm}
\begin{flushleft}
  \color{white}
  {\rmfamily\fontsize{12}{15}\selectfont \coverdoctypename\par}

\vspace{2\baselineskip}
{\sffamily\bfseries\fontsize{18}{22}\selectfont\MakeUppercase{\internshiptitle}\par}

\vspace{5\baselineskip}
{\sffamily\fontsize{12}{15}\selectfont
  MAHASISWA:\par
  \MakeUppercase{\name}\par
  NRP. \nrp\par

  \vspace{3\baselineskip}
  Dosen Pembimbing:\par
  \advisor\par
  NIP. \advisornip\par}

\vspace{4\baselineskip}
{\rmfamily\bfseries\fontsize{14}{17}\selectfont
  Program Studi \studyprogram\par
  Fakultas \faculty\par
  Institut Teknologi Sepuluh Nopember\par
  \the\year\par}

\end{flushleft}
\restoregeometry

\cleardoublepage

% =========================================================================
% HALAMAN SAMPUL DALAM (Mengikuti acuan Tekdok Haura & standar b201)
% =========================================================================
\renewcommand\covercontents{sampul/konten-dalam-magang.tex}
\AddToShipoutPictureBG*{%
  \AtPageLowerLeft{%
    \hspace{45mm}%
    \raisebox{215mm}{\includegraphics[width=40mm]{sampul/gambar/logo-its.png}}%
  }%
}

\pagenumbering{roman}
\setcounter{page}{4}
\thispagestyle{plain}
\newgeometry{top=103mm,left=40mm,right=20mm,bottom=20mm}
  \color{black}
\begin{flushleft}
  \input{\covercontents}
\end{flushleft}
\restoregeometry

\cleardoublepage



\paragraphtextstyle

% =========================================================================
% PENOMORAN HALAMAN BAGIAN AWAL (Sesuai Pedoman ITS & Acuan Haura)
% Dimulai dengan angka Romawi kecil (i, ii, iii, dst.) dari Lembar Pengesahan.
% =========================================================================
\pagenumbering{roman}
\setcounter{page}{1}

% Lembar pengesahan magang sama dengan proposal; judul laporan memakai label khusus.
\thispagestyle{plain}
\phantomsection\addcontentsline{toc}{chapter}{LEMBAR PENGESAHAN}
\begingroup
\paragraphtextstyle
\setlength{\parindent}{0pt}
\setlength{\parskip}{0pt}
\begin{center}
  \textbf{LEMBAR PENGESAHAN\ifdefined\approvaltitlelabel\ LAPORAN\fi}
\end{center}
\ifdefined\approvaltitlelabel
  \vspace{8mm}
  \renewcommand{\arraystretch}{1.8}
\else
  \renewcommand{\arraystretch}{1.5}
\fi
\noindent\begin{tabularx}{\textwidth}{@{}p{0.37\textwidth}cX@{}}
  Mata Kuliah/Kode & : & Magang / KT234602 \\
  \ifdefined\approvaltitlelabel\approvaltitlelabel\else Judul\fi & : & \internshiptitle \\
  Nama Mahasiswa & : & \name \\
  NRP & : & \nrp \\
  Program Studi & : & \studyprogram \\
  Fakultas & : & Fakultas \faculty \\
  Waktu Pelaksanaan Magang & : & \internshipperiod \\
  Tempat Magang & : & \internshipcompany \\
  Alamat Magang & : & \internshipaddress \\
  Telp. & : & \internshipphone \\
  E-mail & : & \internshipemail \\
\end{tabularx}

\ifdefined\approvaltitlelabel
  \vspace{4mm}
  \renewcommand{\arraystretch}{1}
  \noindent\begin{tabularx}{\textwidth}{@{}>{\centering\arraybackslash}X>{\centering\arraybackslash}X@{}}
    & \place, \the\day{} \MONTH{} \the\year \\
    Mengetahui & \\
    Pembimbing Magang & Pembimbing Lapangan \\
    \multicolumn{2}{c}{\rule{0pt}{20mm}} \\
    (\advisor) & (\internshipsupervisor) \\
    NIP. \advisornip & \\
  \end{tabularx}
  \vspace{4mm}
  \begin{center}
    Koordinator Magang\par
    \rule{0pt}{20mm}\par
    (\internshipcoordinator)\par
    NUPTK. \internshipcoordinatornip
  \end{center}
\else
  \vspace{32mm}
  \renewcommand{\arraystretch}{1}
  \noindent\begin{tabularx}{\textwidth}{@{}>{\raggedright\arraybackslash}X>{\raggedleft\arraybackslash}X@{}}
    & \place, \the\day{} \MONTH{} \the\year \\
    Mengetahui & \\
    Koordinator Magang & Dosen Pembimbing Magang \\
    \multicolumn{2}{c}{\rule{0pt}{30mm}} \\
    (\internshipcoordinator) & (\advisor) \\
    NIP. \internshipcoordinatornip & NIP. \advisornip \\
  \end{tabularx}
\fi
\endgroup

\clearpage
% Pernyataan Keaslian Laporan Magang
\phantomsection\addcontentsline{toc}{chapter}{PERNYATAAN KEASLIAN LAPORAN MAGANG}
\begin{center}
  \textbf{PERNYATAAN KEASLIAN LAPORAN MAGANG}
\end{center}

\vspace{4ex}

\noindent Saya mahasiswa Prodi Teknologi Kedokteran, Fakultas Kedokteran dan Kesehatan, Institut Teknologi Sepuluh Nopember, dengan identitas:

\vspace{1ex}

\noindent Nama\quad : \name

\noindent NRP\quad\quad : \nrp

\vspace{4ex}

\noindent Dengan ini menyatakan bahwa laporan magang yang saya buat dengan judul ``\tatitle'' adalah:

\begin{enumerate}[leftmargin=*,label=\arabic*.]
  \item Bukan merupakan duplikasi laporan yang sudah ada, dipublikasikan atau yang pernah dipakai untuk mendapatkan gelar sarjana di universitas lain, kecuali pada bagian-bagian sumber informasi yang dicantumkan sebagai kutipan/referensi dengan cara yang semestinya.
  \item Dibuat dan diselesaikan sendiri dengan menggunakan data-data hasil pelaksanaan magang dalam perusahaan tersebut.
\end{enumerate}

\noindent Demikian pernyataan ini saya buat dan jika terbukti tidak memenuhi persyaratan yang telah saya nyatakan di atas, maka saya bersedia apabila laporan magang ini dibatalkan.

\vspace{8ex}

\begin{flushright}
  \begin{tabular}{c}
    Xxxxxxx, Xx Xxxxxxx 20xx \\
    Yang Membuat Pernyataan \\[8mm]
    TTD + Materai \\[8mm]
    \name \\
    NRP. \nrp
  \end{tabular}
\end{flushright}

\clearpage

% Kata Pengantar Laporan Magang
\begin{center}
  \textbf{KATA PENGANTAR}
\end{center}

\addcontentsline{toc}{chapter}{KATA PENGANTAR}

\vspace{4ex}

\begingroup
\emergencystretch=2em
Kata pengantar tidak boleh lebih dari 1 (satu) halaman, berisi penjelasan tentang maksud penulisan laporan magang, termasuk ucapan terima kasih kepada orang-orang yang punya kontribusi langsung dalam penyusunan laporan magang.
\par
\endgroup

\vspace{16ex}

\begin{flushright}
  \begin{tabular}{c}
    \place{}, \the\day{} \MONTH{} \the\year{} \\
    \\
    \\
    \\
    \\
    \\
    \\
    Penulis
  \end{tabular}
\end{flushright}

\clearpage



% Daftar isi, gambar, tabel, dan simbol
% Daftar awal memakai Heading 1 untuk judul dan 11 pt untuk isi.
\renewcommand*\contentsname{DAFTAR ISI}
\renewcommand*\listfigurename{DAFTAR GAMBAR}
\renewcommand*\listtablename{DAFTAR TABEL}

\begingroup
\fontsize{11}{13}\selectfont
\setlength{\cftaftertoctitleskip}{79pt}
\renewcommand{\cfttoctitlefont}{\hfill\normalfont\bfseries\fontsize{12}{15}\selectfont\color{black}}
\renewcommand{\cftaftertoctitle}{\hfill}
\renewcommand{\cftchapleader}{\cftdotfill{\cftdotsep}}
\renewcommand{\cftsecleader}{\cftdotfill{\cftdotsep}}
\renewcommand{\cftsubsecleader}{\cftdotfill{\cftdotsep}}
\renewcommand{\cftsubsubsecleader}{\cftdotfill{\cftdotsep}}
\phantomsection\addcontentsline{toc}{chapter}{\contentsname}
\tableofcontents
\endgroup
\cleardoublepage

\begingroup
\fontsize{11}{13}\selectfont
\phantomsection\addcontentsline{toc}{chapter}{\listfigurename}
\listoffigures
\endgroup
\cleardoublepage

\begingroup
\fontsize{11}{13}\selectfont
\phantomsection\addcontentsline{toc}{chapter}{\listtablename}
\listoftables
\endgroup
\cleardoublepage

\begingroup
\fontsize{11}{13}\selectfont
\chapter*{DAFTAR SIMBOL}
\phantomsection\addcontentsline{toc}{chapter}{DAFTAR SIMBOL}
\noindent\begin{tabular}{@{}ll@{}}
  $\mathbf{F}$ & Gaya \\[1ex]
  $\mathbf{v}$ & Kecepatan \\[1ex]
  $t$ & Waktu
\end{tabular}
\endgroup
\cleardoublepage


% Nomor halaman isi dimulai dari angka Arab 1
\pagenumbering{arabic}

% Bab 1 Pendahuluan
\chapter{PENDAHULUAN}
\label{chap:pendahuluan}

\paragraphtextstyle

\section{LATAR BELAKANG}
\label{sec:latarbelakang}

Ditulis latar belakang laporan magang. Latar belakang menjelaskan permasalahan, mengapa masalah tersebut penting untuk dipecahkan dan lain-lain.

\section{PERMASALAHAN}
\label{sec:permasalahan}

Permasalah yang diuraikan berupa permasalahan yang dihadapi selama proses magang di tempat Anda.

\section{TUJUAN}
\label{sec:tujuan}

Tujuan diuraikan dengan contoh sebagai berikut:
\begin{enumerate}[label=\alph*.]
  \item Memberikan gambaran dunia kerja yang sebenarnya kepada mahasiswa sebagai bekal untuk kemudian hari.
\end{enumerate}

\section{MANFAAT MAGANG}
\label{sec:manfaat}

\subsection{Bagi Mahasiswa}
\begin{enumerate}[label=\alph*.]
  \item Dapat memperoleh kesempatan untuk menerapkan ilmu pengetahuan yang telah diperoleh dalam perkuliahan untuk diterapkan dalam lapangan kerja.
  \item dan seterusnya
\end{enumerate}

\subsection{Bagi Perguruan Tinggi}
\begin{enumerate}[label=\alph*.]
  \item Mengenalkan dunia profesional terkait alat kesehatan kepada mahasiswa agar tercipta mahasiswa yang unggul, berpengalaman dan berkualitas sehingga siap terjun ke dunia kerja.
  \item dan seterusnya
\end{enumerate}

\subsection{Bagi Perusahaan}
\begin{enumerate}[label=\alph*.]
  \item Mahasiswa dapat membantu pekerjaan/proyek.
  \item dan seterusnya
\end{enumerate}

\clearpage

% Bab 2 Tinjauan Pustaka
%! TeX root = ../main.tex
\chapter{PROFIL PERUSAHAAN}
\label{chap:profilperusahaan}

\paragraphtextstyle

\section{IDENTITAS PERUSAHAAN}
\label{sec:identitasperusahaan}

Nama perusahaan alat kesehatan tempat penulis melakukan magang adalah di \ldots. Alamat perusahaan alat kesehatan tempat penulis melakukan magang adalah di \ldots.

\section{SEJARAH PERUSAHAAN}
\label{sec:sejarahperusahaan}

XXX adalah perusahaan XXXXXX yang berdiri pada tahun XXXXX. Perusahaan ini didirikan oleh \ldots dan seterusnya.

\section{BIDANG USAHA}
\label{sec:bidangusaha}

XXX secara khusus bergerak di bidang XXXXX.

\section{SUSUNAN ORGANISASI}
\label{sec:susunanorganisasi}

\begin{figure}[H]
  \centering
  \fbox{\parbox[c][45mm][c]{0.85\textwidth}{\centering Struktur Organisasi Perusahaan}}
  \caption{Struktur Organisasi PT Mega Andalan Kalasan}
  \label{fig:strukturorganisasi}
  {\fontsize{10}{12}\selectfont Sumber: Kepala Unit Bagian HRD (xxxxx, 20xx)\par}
\end{figure}

\section{PORTOFOLIO PERUSAHAAN}
\label{sec:portofolioperusahaan}

Perusahaan yang telah berdiri selama xx tahun ini memiliki beragam produk yang telah dihasilkan. Berikut beberapa proyek yang pernah ditangani oleh Artmosphere Interior Design Inc.

\clearpage
\begin{table}[H]
  \centering
  \caption{Portofolio PT Mega Andalan Kalasan}
  \label{tab:portofolio}
  \begin{tabularx}{\textwidth}{|c|c|X|}
    \hline
    No & Tahun & Nama Proyek \\
    \hline
    1 & 2016 & Tempat Tidur Rumah Sakit \\
    \hline
    2 & 2016 & Meja Operasi \\
    \hline
    3 & 2016 & Alat Transportasi Pasien \\
    \hline
  \end{tabularx}

  \vspace{0.5ex}
  {\fontsize{10}{12}\selectfont Sumber: Kepala Unit Bagian HRD (xxxxx, 20xx)\par}
\end{table}

\clearpage

% Bab 3 Metodologi
%! TeX root = ../main.tex
\chapter{LANGKAH DAN URAIAN KEGIATAN}
\label{chap:langkahuraian}

\paragraphtextstyle

\section{BIDANG KERJA}
\label{sec:bidangkerja}

\section{LANGKAH DAN URAIAN KEGIATAN}
\label{sec:langkahuraian}

\section{JADWAL KEGIATAN}
\label{sec:jadwalkegiatan}

\clearpage

% Bab 4 Hasil dan Pembahasan
\chapter{HASIL DAN PEMBAHASAN}
\label{chap:hasildanpembahasan}

\paragraphtextstyle

\section{HASIL}
\label{sec:hasil}

\section{PEMBAHASAN}
\label{sec:pembahasan}

\clearpage

% Bab 5 Kesimpulan dan Saran
\chapter{KESIMPULAN DAN SARAN}
\label{chap:kesimpulansaran}

\paragraphtextstyle
\renewcommand{\thesection}{4.\arabic{section}}

\section{KESIMPULAN}
\label{sec:kesimpulan}

\section{SARAN}
\label{sec:saran}

\clearpage

% Daftar Pustaka
\chapter*{DAFTAR PUSTAKA}
\addcontentsline{toc}{chapter}{DAFTAR PUSTAKA}
\renewcommand\refname{}
\vspace{2ex}
\renewcommand{\bibname}{}
\begingroup
\def\chapter*#1{}
\paragraphtextstyle
\nocite{rajpurkar2022ai}% Contoh daftar pustaka APA; hapus setelah ada sitasi sendiri.
\printbibliography
\endgroup
\clearpage

% Lampiran laporan magang
\begin{center}
  \textbf{LAMPIRAN}
\end{center}

\addcontentsline{toc}{chapter}{LAMPIRAN}

\vspace{1.5ex}

\begin{enumerate}[topsep=0pt]
  \item Surat Penerimaan Magang
  \item Surat Keterangan Telah Menyelesaikan Magang
  \item Surat Pernyataan Tidak Melanggar Etika
  \item Dokumentasi saat Magang
\end{enumerate}

\clearpage

% Biodata Penulis
\begingroup
\def\cvtitle{BIODATA PENULIS}
%! TeX root = ../proposal_magang.tex
\providecommand{\cvtitle}{CURRICULUM VITAE}
\phantomsection
\addcontentsline{toc}{chapter}{\cvtitle}
\noindent\makebox[\textwidth][c]{\textbf{\cvtitle}}\par
\vspace{1ex}

\begingroup
\paragraphtextstyle
\fontsize{12}{13.5}\selectfont
\setlength{\parindent}{0pt}
\setlength{\parskip}{0pt}
\noindent\begin{tabularx}{\textwidth}{@{}>{\raggedright\arraybackslash}p{42mm}>{\raggedright\arraybackslash}X@{}}
\multicolumn{2}{@{}l@{}}{\textbf{Data Diri}} \\
Nama Lengkap & [Nama Lengkap] \\
Alamat & [Alamat Lengkap] \\
Nomor Telepon & [Nomor Telepon] \\
Email & [Email] \\
LinkedIn/Portofolio & [Tautan LinkedIn atau Portofolio (jika ada)] \\
\multicolumn{2}{@{}l@{}}{} \\[3pt]
\multicolumn{2}{@{}l@{}}{\textbf{Profil}} \\
& [Tuliskan ringkasan singkat tentang diri Anda, misalnya: ``Saya adalah mahasiswa [jurusan] di [universitas] dengan minat dalam [bidang]. Saya mencari kesempatan magang untuk mengembangkan keterampilan dan pengalaman di bidang tersebut.''] \\
\multicolumn{2}{@{}l@{}}{} \\[3pt]
\multicolumn{2}{@{}l@{}}{\textbf{Pendidikan}} \\
Perguruan Tinggi & Institut Teknologi Sepuluh Nopember \\
Program Studi & Teknologi Kedokteran \\
Tahun Studi & [Tahun Masuk -- Tahun Sekarang] \\
IPK & [Nilai IPK] \\
\multicolumn{2}{@{}l@{}}{} \\[3pt]
\multicolumn{2}{@{}l@{}}{\textbf{Pengalaman Kerja [dihapus jika tidak ada]}} \\
Nama Organisasi & [Nama Perusahaan] \\
Jabatan & [Jabatan] \\
Periode & [Bulan, Tahun Mulai -- Bulan, Tahun Selesai] \\
Deskripsi Tugas/Pencapaian & [Deskripsikan tugas utama dan pencapaian relevan] \\
\multicolumn{2}{@{}l@{}}{} \\[3pt]
\multicolumn{2}{@{}l@{}}{\textbf{Keterampilan}} \\
Keterampilan Teknis & - [Sebutkan keterampilan teknis seperti perangkat lunak, coding, dll.] \\
Keterampilan Interpersonal & - [Sebutkan keterampilan seperti komunikasi, kerja tim, dll.] \\
\multicolumn{2}{@{}l@{}}{} \\[3pt]
\multicolumn{2}{@{}l@{}}{\textbf{Pengalaman Organisasi dan Ekstrakurikuler}} \\
Nama Organisasi & [Nama Organisasi atau Komunitas yang Diikuti] \\
Jabatan/Peran & [Jabatan atau Peran dalam organisasi tersebut] \\
Periode Keaktifan & [Tahun Mulai -- Tahun Selesai/Sekarang] \\
Deskripsi Kontribusi & - [Deskripsikan peran dan kontribusi] \\
\multicolumn{2}{@{}l@{}}{} \\[3pt]
\multicolumn{2}{@{}l@{}}{\textbf{Prestasi [dihapus jika tidak ada]}} \\
\multicolumn{2}{@{}p{\textwidth}@{}}{- [Nama kegiatan, penyelenggara, tingkat, waktu pelaksanaan, pencapaian]}
\end{tabularx}
\endgroup

\endgroup

\clearpage

\end{document}
